# Abstract  
## Topological Integrity Gravity (TIG)

---

## 🇬🇧 English Version

We introduce Topological Integrity Gravity (TIG), a spectrally derived effective extension of General Relativity in which curvature corrections emerge from the spectral structure of the Dirac operator and are constrained by an intrinsic stability functional.

In the low-curvature limit, TIG reduces to a curvature-corrected \(f(R)\)-type theory and reproduces the structure of Starobinsky gravity. However, in contrast to phenomenological \(f(R)\) models, the form of the corrections in TIG is not freely chosen but follows from a spectral construction augmented by an integrity functional of the form
\[
\mathcal{I}[D] = \mathrm{Tr}\log\left(1 + \frac{D^2}{\Lambda^2}\right),
\]
which introduces a non-polynomial, scale-dependent feedback mechanism.

This structure leads to a stability-constrained parameter space and predicts a scalar curvature degree of freedom with mass
\[
m = \frac{\Lambda}{\sqrt{6\alpha}},
\]
whose coupling to matter is determined by the same parameter combination, reducing model freedom and enabling direct observational tests.

We analyze a class of static, spherically symmetric solutions exhibiting a regular de Sitter core, multiple horizons, and an extremal configuration with vanishing Hawking temperature. This implies the existence of a remnant endpoint for black hole evaporation.

While the resulting solutions are shown to be kinematically regular and thermodynamically consistent, a full dynamical stability analysis remains an open problem. We outline a perturbative framework for future investigation.

TIG therefore provides a spectrally motivated selection principle for curvature corrections beyond phenomenological \(f(R)\) gravity, linking geometric structure, stability, and observational signatures within a unified effective framework.

---

## 🇩🇪 Deutsche Version

Wir stellen Topological Integrity Gravity (TIG) vor, eine spektral hergeleitete effektive Erweiterung der Allgemeinen Relativitätstheorie, in der Krümmungskorrekturen aus der spektralen Struktur des Dirac-Operators entstehen und durch ein intrinsisches Integritätsfunktional eingeschränkt werden.

Im Niedrigkrümmungs-Grenzfall reduziert sich TIG auf eine \(f(R)\)-artige Theorie und reproduziert die Struktur der Starobinsky-Gravitation. Im Gegensatz zu phänomenologischen \(f(R)\)-Modellen sind die Korrekturterme jedoch nicht frei wählbar, sondern folgen aus einer spektralen Konstruktion, ergänzt durch ein Funktional der Form
\[
\mathcal{I}[D] = \mathrm{Tr}\log\left(1 + \frac{D^2}{\Lambda^2}\right),
\]
das eine nicht-polynomiale, skalenabhängige Rückkopplung einführt.

Diese Struktur führt zu einem stabilitätsbeschränkten Parameterraum und sagt einen skalaren Krümmungsmodus mit Masse
\[
m = \frac{\Lambda}{\sqrt{6\alpha}}
\]
vorher, dessen Kopplung an Materie durch dieselbe Parameterkombination bestimmt wird. Dadurch wird die Modellfreiheit reduziert und eine direkte experimentelle Testbarkeit ermöglicht.

Wir analysieren eine Klasse statischer, sphärisch symmetrischer Lösungen mit regulärem de-Sitter-Kern, mehreren Horizonten und einem extremalen Zustand mit verschwindender Hawking-Temperatur. Dies impliziert die Existenz eines Remnant-Endzustands für die Schwarze-Loch-Verdampfung.

Während die Lösungen kinematisch regulär und thermodynamisch konsistent sind, bleibt eine vollständige dynamische Stabilitätsanalyse eine offene Aufgabe. Ein entsprechender perturbativer Zugang wird skizziert.

TIG liefert damit ein spektral motiviertes Selektionsprinzip für Krümmungskorrekturen über phänomenologische \(f(R)\)-Modelle hinaus und verbindet geometrische Struktur, Stabilität und beobachtbare Signaturen in einem einheitlichen effektiven Rahmen.
